\section{Integración del Prototipo (Objetivo Específico 3)}

\subsection{Niveles de Integración y Ensamblaje}
La integración del Rover Olympus se rige por un enfoque de "V-Model", donde cada nivel de ensamblaje debe ser verificado antes de pasar al siguiente.

\begin{enumerate}
    \item \textbf{Nivel 1 (Componentes):} Verificación individual de actuadores, sensores y módulos de potencia.
    \item \textbf{Nivel 2 (Subsistemas):} Integración del bus de comunicaciones CSP/CAN y validación de la telemetría base entre el microcontrolador y el C\&DH.
    \item \textbf{Nivel 3 (Sistema Completo):} Montaje final de la electrónica sobre el chasis y carga del firmware/software de navegación.
\end{enumerate}

\subsection{Entorno de Integración de Software (SITL)}
Para reducir riesgos de colisión en el prototipo físico, se utiliza un entorno de simulación \textbf{Software-in-the-Loop (SITL)} que permite validar la lógica de control sin riesgo para el hardware:

\begin{itemize}
    \item \textbf{Simulador:} Gazebo / Isaac Sim con modelos físicos que emulan la baja gravedad y la fricción del regolito.
    \item \textbf{Gemelo Digital:} Modelo URDF (Unified Robot Description Format) que refleja con precisión las dimensiones, centros de masa y límites de articulación definidos en la arquitectura física.
    \item \textbf{Validación Pre-Hardware:} Los algoritmos de Nav2 y SLAM-Toolbox deben superar escenarios de autonomía en el simulador (identificación de obstáculos y planificación de rutas) antes de su implementación en el hardware físico.
\end{itemize}

\subsection{Protocolo de Integración de Hardware}
Para asegurar la integridad del rover durante el ensamblaje final, se sigue un protocolo estricto de energización y verificación por etapas:

\begin{table}[H]
    \centering
    \begin{tabularx}{\textwidth}{|l|X|l|}
        \hline
        \rowcolor{AgencyBlue!10} \textbf{Etapa} & \textbf{Actividad Crítica} & \textbf{Requerimiento Validado} \\ \hline
        \textbf{Mecánica} & Ensamble de chasis y transmisión. & RF-MOV-01 \\ \hline
        \textbf{Potencia} & Verificación de voltajes y BMS sin carga. & RF-ENE-01 \\ \hline
        \textbf{Cómputo} & Boot de Linux (Ubuntu 22.04 LTS) y ROS 2. & RI-SW-01 \\ \hline
        \textbf{Navegación} & Calibración de LiDAR e IMU sobre el rover. & RF-GNC-03 \\ \hline
        \textbf{Comunicación} & Telemetría bidireccional mediante CSP. & RI-USR-01 \\ \hline
    \end{tabularx}
    \caption{Secuencia de integración y validación progresiva.}
\end{table}
